% Generated draft: requires mathematical and source review.
\documentclass[11pt]{article}
\usepackage[margin=25mm]{geometry}
\usepackage{fontspec}
\setmainfont{Latin Modern Roman}
\usepackage{amsmath,amssymb,mathtools,unicode-math}
\setmathfont{Latin Modern Math}
\usepackage{enumitem,needspace,xurl,hyperref}
\setlength{\parindent}{0pt}
\setlength{\parskip}{5pt}
\setlength{\emergencystretch}{3em}
\newcommand{\N}{\mathbb{N}}
\newcommand{\Z}{\mathbb{Z}}
\newcommand{\Q}{\mathbb{Q}}
\newcommand{\R}{\mathbb{R}}
\newcommand{\C}{\mathbb{C}}
\newcommand{\F}{\mathbb{F}}
\newcommand{\A}{\mathbb{A}}
\newcommand{\PP}{\mathbb{P}}
\newcommand{\E}{\mathbb{E}}
\newcommand{\eps}{\varepsilon}
\newcommand{\dd}{\,\mathrm{d}}
\begin{document}
\section*{UnsolvedMath ID: 30005410}
\textbf{Source problem code:} catalogue entry
\subsection{Definitions and mathematical statement}
For a commutative Noetherian ring \(R\), let \(\operatorname{mod}R\) denote the category of finitely generated \(R\)-modules, let \(D^b(\operatorname{mod}R)\) be its bounded derived category, and let \(\operatorname{Perf}(R)\) be its full subcategory of complexes quasi-isomorphic to bounded complexes of finitely generated projective modules. The singularity category is the Verdier quotient
\[
D_{\mathrm{sg}}(R)=D^b(\operatorname{mod}R)/\operatorname{Perf}(R).
\]
Two rings are called singular equivalent when their singularity categories are equivalent as triangulated categories. The supplied excerpt does not specify whether equivalences must additionally be \(\C\)-linear or preserve an enhancement; those requirements remain unresolved.

\textbf{Question 3 (supplied source).} How fine is the invariant \(R\mapsto D_{\mathrm{sg}}(R)\)? Determine when two rings \(R,S\) satisfy
\[
D_{\mathrm{sg}}(R)\simeq D_{\mathrm{sg}}(S).
\]
The source's footnote restricts its displayed list of known examples to nonisomorphic commutative complete local Noetherian \(\C\)-algebras. This is therefore the principal setting indicated for the classification question, although the question itself says only ``rings.'' In this setting, the target is to describe the equivalence relation on such rings induced by their singularity categories, including equivalences with \(R\not\cong S\). Allowing noncommutative rings is discussed separately below; the excerpt does not formulate a precise classification target for that broader setting.

\textbf{Source notation and examples.} For a positive integer \(n\), a primitive \(n\)-th root of unity \(\epsilon_n\in\C\), and \((a_1,\ldots,a_m)\in\Z_{>0}^m\), write
\[
G_n(a_1,\ldots,a_m)=\left\langle\operatorname{diag}(\epsilon_n^{a_1},\ldots,\epsilon_n^{a_m})\right\rangle\subseteq\operatorname{GL}(m,\C).
\]
Its diagonal action on \(\C[[z_1,\ldots,z_m]]\) gives the invariant ring \(\C[[z_1,\ldots,z_m]]^{G_n(a_1,\ldots,a_m)}\). The source denotes this group by \(\frac1n(a_1,\ldots,a_m)\) and calls it a group of order \(n\). For the stated unrestricted tuples, its order is actually \(n/\gcd(n,a_1,\ldots,a_m)\); the order assertion needs the additional condition \(\gcd(n,a_1,\ldots,a_m)=1\).

The source lists the following examples, as historical context rather than a classification:
\begin{itemize}
\item[(E0)] \(D_{\mathrm{sg}}(R)\simeq0\simeq D_{\mathrm{sg}}(S)\) for regular rings \(R,S\), attributed to Auslander--Buchsbaum and Serre. The excerpt does not state the full ambient hypotheses for this contextual assertion.
\item[(E1)] The equivalence attributed to Kn\"orrer (1987), reference [10], is
\[
D_{\mathrm{sg}}\!\left(\frac{\C[[z_0,\ldots,z_d]]}{(f)}\right)
\simeq
D_{\mathrm{sg}}\!\left(\frac{\C[[z_0,\ldots,z_d,x_1,\ldots,x_{2m}]]}{(f+x_1^2+\cdots+x_{2m}^2)}\right),
\qquad 0\ne f\in\C[[z_0,\ldots,z_d]].
\]
The supplied excerpt does not explicitly give the ranges of \(d,m\) or require \(f\) to be a nonunit. These omissions are retained as qualifications rather than silently strengthened hypotheses.
\item[(E2)] The extracted display is read as
\[
D_{\mathrm{sg}}\!\left(\C[[y_1,y_2]]^{G_n(1,1)}\right)
\simeq
D_{\mathrm{sg}}\!\left(\frac{\C[z_1,\ldots,z_{n-1}]}{(z_1,\ldots,z_{n-1})^2}\right).
\]
It is attributed to Yang (2015), reference [12], Kawamata (2015), reference [9], and Kalck--Karmazyn (2017), reference [6]. The source says Yang's result is based on joint works [7,8]. A further footnote concerning a generalization is truncated. The range of \(n\), the display's transcription, and the scope of that generalization require checking against the report.
\item[(E3)] The source attributes to Kalck (2021), reference [3], the equivalence
\[
D_{\mathrm{sg}}\!\left(\C[[x_1,x_2,x_3]]^{G_2(1,1,1)}\right)
\simeq
D_{\mathrm{sg}}\!\left(\C[[y_1,y_2]]^{G_4(1,1)}\right).
\]
The source remarks that this equivalence, in contrast to (E1) and (E2), does not preserve the parity of Krull dimension.
\end{itemize}
The source's description of its list as complete is qualified by its footnote: it concerns singular equivalences between nonisomorphic commutative complete local Noetherian \(\C\)-algebras that the source authors could find in the literature. It also points to an older list in [11, Example 4.1]. Neither description supplies a classification theorem.

The source further discusses proofs of (E2) and (E3) using noncommutative resolutions and relative singularity categories, building on [7,8]. It alternatively describes maximal Cohen--Macaulay modules \(M_R,M_S\) as unique up to taking syzygies and direct summands, with \(M_R,M_S\) direct summands of \(\Omega(\omega_R),\Omega(\omega_S)\), respectively, where \(\omega_R,\omega_S\) are canonical modules and \(\Omega\) denotes a syzygy. The exact scope and meaning of this uniqueness formulation are not specified sufficiently in the excerpt to make it a separate precise assertion.

If \(S\) is allowed to be noncommutative, the source reports that for every ADE surface singularity \(R\) except type \(E_8\), its cited work [5] constructs an uncountable family of complete Noetherian \(\C\)-algebras \(S_R\) singular equivalent to \(R\), with pairwise nonisomorphic centers. The necessary conventions for singularity categories in this broader setting and the full reference [5] are not supplied.

The opening fragment about two-periodic stable ranges induced by matrix factorizations of \(f\) refers to an absent condition (P2); it does not provide an additional well-specified hypothesis for Question 3. The registry labels the problem open and supplies a dated assessment that the examples and subsequent constructions leave the general commutative classification unresolved. These are supplied curation claims, not independently verified mathematical or literature conclusions.
\subsection{Short English statement}
Classify commutative complete local Noetherian \(\C\)-algebras according to equivalence of their singularity categories, which measure the part of the bounded derived category remaining after perfect complexes are removed. The source gives examples of nonisomorphic rings with equivalent singularity categories and discusses a broader noncommutative setting. Its exact equivalence conventions, some example hypotheses, and several bibliography details remain unresolved. The supplied open status requires independent review.
\subsection{Sources}
\begin{itemize}
\item Local registry: \texttt{data/unsolvedmath/problems.json}, permanent UnsolvedMath ID \(30005410\), source code \texttt{catalogue entry}, title ``Classifying Rings Up to Singularity Equivalence.'' The supplied record gives proposed year 2023 and no named proposer.
\item Supplied report citation: \emph{Resolutions in Local Algebra and Singularity Theory} (2023), Oberwolfach Reports, pp.~327--395, DOI 10.4171/owr/2023/6. \url{https://doi.org/10.4171/owr/2023/6}. The extracted passage contains Question 3 and a page marker for p.~349.
\item Supplied literature source: ``OWR-2023-6, Question 3.'' \url{https://oa.tib.eu/renate/server/api/core/bitstreams/11773ea5-b1e5-4e9e-b650-808e9d9dc944/content}.
\item Supplied literature source: ``Kalck, new singular equivalence.'' \url{https://arxiv.org/abs/2103.06584}.
\item Attributions retained from the excerpt: Auslander--Buchsbaum and Serre; Kn\"orrer (1987), [10]; Yang (2015), [12]; Kawamata (2015), [9]; Kalck--Karmazyn (2017), [6]; and Kalck (2021), [3]. References [5], [7], [8], and [11, Example 4.1] also occur. Their full bibliography entries are not supplied and have not been reconstructed.
\item Literature assessment dated exactly \texttt{2026-08-22}: the supplied curation labels the problem open and says the displayed examples and subsequent singular-equivalence constructions leave the general commutative classification unresolved. This is supplied dated curation, not independently verified. The background explicitly requests independent verification of time-sensitive status and cited claims.
\item The registry labels the statement \texttt{corrected\_verified} and reports that the official report was checked directly. This records the supplied verification claim; this draft does not claim to have visited or independently checked any source.
\end{itemize}
\end{document}
